\chapter{Le desk d'ajustement de valeur}\label{ch:rc:the-valuation-adjustment-desk}

Un commercial chiffre un swap de vingt ans pour une entreprise industrielle non notée qui
souhaite fixer le taux d'un long emprunt. Le desk de taux cote le taux au pair en quelques secondes. Puis le
desk d'ajustement de valeur ajoute sa charge~: le risque de crédit de l'entreprise, le financement de la
position non collatéralisée par la banque, le capital que la transaction immobilisera pendant deux décennies. Dans
l'exemple de ce chapitre, la charge atteint près de quarante points de base par an, davantage que
la fourchette acheteur-vendeur du swap lui-même, et le client peut se tourner vers une banque qui a une vision plus
favorable de l'un de ces trois éléments. Les chapitres~17 à~19 ont construit les ajustements~; ce chapitre porte sur le desk
qui en a la charge~: pourquoi il est central, comment il tarifie une nouvelle transaction par rapport à tout ce que la banque détient déjà
avec le client, ce qu'il peut couvrir et ce qu'il ne peut pas couvrir, et comment ses charges parviennent aux
desks de trading.

\section{Pourquoi un desk central}

\begin{definition}[Desk XVA]\label{def:rc:the-valuation-adjustment-desk:desk}
Un \emph{desk XVA}\index{desk XVA} est le desk de trading qui détient les ajustements de
valeur d'une banque~: il facture à chaque nouvelle transaction ses ajustements, et porte
de financement et de capital qui en résultent sur son livre, couvre ce qui peut l'être, et déclare les
variations de ces ajustements comme son P\&L.
\end{definition}

Les ajustements sont des propriétés des \omterm{def:rc:counterparty-exposure:netting}{ensembles de compensation}, et non des transactions~: un swap vendu par le desk de taux
et un forward de change vendu par le desk de change au même client se compensent mutuellement, partagent le
collatéral du client et son défaut. Seul un desk qui voit l'\omterm{def:rc:counterparty-exposure:netting}{ensemble de compensation} dans son intégralité peut chiffrer
l'effet d'une transaction sur celui-ci. La centralisation concentre aussi l'expertise et la couverture, qui
exige des credit default swaps, une protection sur indices et des couvertures de marché dans toutes les classes d'actifs
(la \cref{fig:rc:the-valuation-adjustment-desk:org}).

\begin{omfigure}
\begin{tikzpicture}[font=\small,box/.style={draw,rounded corners,minimum width=2.3cm,minimum height=0.85cm,align=center}]
\node[box] (r) at (0,1.2) {desk de taux};
\node[box] (f) at (0,0) {desk de change};
\node[box] (e) at (0,-1.2) {autres desks};
\node[box,fill=omDef!10,minimum height=1.6cm] (x) at (5,0) {desk XVA\\(ensembles de compensation,\\CVA, FVA, KVA)};
\node[box] (m) at (10,1.2) {marché du crédit\\(CDS, indice)};
\node[box] (h) at (10,0) {couvertures de taux\\et de change};
\node[box,fill=omThm!10] (t) at (10,-1.2) {trésorerie et\\capital};
\foreach \n in {r,f,e} {\draw[->,thick] (\n.east) -- (x.west);}
\node[above,align=center,font=\footnotesize] at (2.5,1.05) {transactions, et paie\\la charge};
\draw[->,thick] (x.east) -- (m.west);
\draw[->,thick] (x.east) -- (h.west);
\draw[<->,thick] (x.east) -- (t.west);
\node[above,font=\footnotesize] at (7.5,1.05) {couvertures};
\end{tikzpicture}

\omcaption{Le \omterm{def:rc:the-valuation-adjustment-desk:desk}{desk XVA} entre les desks de trading et les marchés~: les desks lui transmettent
les risques de contrepartie, de financement et de capital de leurs transactions et lui versent une charge~; il couvre
le risque de crédit et le risque de marché sur le marché et règle le financement et le capital avec la trésorerie.
Schéma.}\label{fig:rc:the-valuation-adjustment-desk:org}
\end{omfigure}

\section{Tarification marginale et allocation}

\begin{definition}[XVA incrémentale]\label{def:rc:the-valuation-adjustment-desk:incremental}
La \emph{XVA incrémentale}\index{XVA incrémentale} d'une nouvelle transaction est la variation
des ajustements de l'\omterm{def:rc:counterparty-exposure:netting}{ensemble de compensation} (ou du portefeuille) lorsque la transaction est ajoutée~:
$\mathrm{XVA}(\text{ensemble}+\text{transaction})-\mathrm{XVA}(\text{ensemble})$. C'est ce que la transaction coûte à la banque,
et elle peut être négative lorsqu'elle compense une exposition existante.
\end{definition}

\begin{example}[Deux nouveaux swaps dans un même ensemble de compensation]\label{ex:rc:the-valuation-adjustment-desk:incremental}
L'\omterm{def:rc:counterparty-exposure:netting}{ensemble de compensation} du chapitre~17, avec la contrepartie du chapitre~18, a une CVA de USD~891\,342. Un nouveau
swap de cinq ans de USD~50 millions, sur lequel la banque paie le taux fixe, a une CVA autonome de 37\,061,
mais son ajout ramène la CVA de l'ensemble à 869\,582~: sa CVA incrémentale vaut $-\text{USD}~21\,761$, une
ristourne. Le même swap, la banque recevant le taux fixe, a une CVA autonome de 30\,742 et une
CVA incrémentale de $+22\,784$.
\end{example}

\omcode{code/firm/xvaquote/firm_xvaquote.py}{39}{50}{Allocation d'Euler de la CVA d'un ensemble de compensation à ses transactions.}\label{lst:rc:the-valuation-adjustment-desk:euler}

Les charges marginales dépendent de l'ordre d'arrivée~: la première transaction avec un client acquitte la
charge autonome complète, tandis qu'une transaction compensatrice ultérieure dégage une ristourne. Pour déclarer la
part de chaque transaction dans un ajustement existant, le desk a besoin d'une allocation dont la somme est exacte.

\begin{definition}[Allocation d'Euler]\label{def:rc:the-valuation-adjustment-desk:euler}
L'\emph{allocation d'Euler}\index{allocation d'Euler} attribue à la transaction $i$ la
contribution $a_i = \frac{d}{d\epsilon}\mathrm{XVA}(\ldots,(1+\epsilon)V_i,\ldots)|_{\epsilon=0}$. Lorsque l'ajustement est
homogène de degré un par rapport aux tailles des transactions, comme la CVA l'est, le théorème d'Euler fait que les
contributions s'additionnent au total~; pour la CVA,
$a_i = (1-R)\sum_k\E[D(0,t_k)V_i(t_k)\mathbf 1_{V(t_k)>0}]\,\Delta\mathrm{PD}_k$.
\end{definition}

\begin{example}[Allocation de l'ensemble de compensation]\label{ex:rc:the-valuation-adjustment-desk:euler}
Sur les USD~891\,342 de CVA de l'ensemble, l'\omterm{def:rc:the-valuation-adjustment-desk:euler}{allocation d'Euler} attribue 160\,167 au swap et 731\,175 au
swap de devises, contre des CVA autonomes de 288\,250 et 768\,334
(la \cref{fig:rc:the-valuation-adjustment-desk:euler}). Le swap, qui compense une partie de
l'exposition en devises, reçoit l'essentiel du bénéfice de compensation.
\end{example}

\begin{omfigure}
\begin{tikzpicture}
\begin{axis}[width=11.5cm,height=5.2cm,ybar,bar width=24pt,
  symbolic x coords={swap,cross-currency},xtick=data,xticklabels={swap,{swap de devises}},enlarge x limits=0.5,
  ylabel={CVA (milliers d'USD)},
  tick label style={font=\small},label style={font=\small},
  ymin=0,ymax=900,grid=major,grid style={gray!25},
  nodes near coords,every node near coord/.append style={font=\footnotesize,/pgf/number format/fixed,/pgf/number format/precision=0},
  legend columns=2,legend style={font=\footnotesize,at={(0.5,-0.2)},anchor=north,draw=none,fill=none,/tikz/every even column/.append style={column sep=8pt}},
  table/col sep=comma]
\addplot[fill=omExo!40,draw=omExo] table[x=trade,y=standalone]{figdata/rates-credit-risk/20-the-valuation-adjustment-desk/euler.csv};
\addplot[fill=omDef!60,draw=omDef] table[x=trade,y=euler]{figdata/rates-credit-risk/20-the-valuation-adjustment-desk/euler.csv};
\legend{autonome,contribution d'Euler}
\end{axis}
\end{tikzpicture}

\omcaption{CVA autonome de chaque transaction comparée à sa contribution d'Euler à la CVA de
l'\omterm{def:rc:counterparty-exposure:netting}{ensemble de compensation}. Les contributions s'additionnent à la CVA de l'ensemble~; l'écart avec les
chiffres autonomes est le bénéfice de compensation, réparti proportionnellement au rôle de chaque transaction dans l'exposition. Données~:
le tutoriel du chapitre.}\label{fig:rc:the-valuation-adjustment-desk:euler}
\end{omfigure}

\section{Ce que la CVA couvre et ce qu'elle ne couvre pas}

La CVA comporte deux types de risque~: les facteurs de marché qui font varier l'exposition, couverts par des swaps,
options et de change, comme tout livre de dérivés, et le crédit de la contrepartie, couvert par des
instruments de crédit. Beaucoup de contreparties n'ont pas de CDS négociés~: le desk a besoin d'un spread pour les
évaluer et d'un proxy pour les couvrir.

\begin{definition}[Spread de crédit de substitution]\label{def:rc:the-valuation-adjustment-desk:proxy}
Un \emph{spread de crédit de substitution}\index{spread de crédit de substitution} (proxy credit spread) est un spread attribué à une
contrepartie sans cotations de crédit liquides, estimé à partir de pairs liquides qui lui ressemblent en
qualité de crédit, secteur et zone géographique, ou rattaché à un seul nom apparenté. Dans l'Union européenne, un proxy
utilisé pour le capital CVA doit tenir compte des trois critères, avec au moins trois catégories sectorielles (public,
financier, autre) et quatre zones géographiques (Europe, Amérique du Nord, Asie, reste du monde).
\end{definition}

\begin{example}[Un proxy pour le client non noté]\label{ex:rc:the-valuation-adjustment-desk:proxy}
On régresse le logarithme de dix-huit spreads à dix ans de pairs illustratifs sur les indicateurs de notation (A, BBB, BB),
de secteur (industriels, services publics, consommation) et de zone (États-Unis, UE). L'erreur de l'ajustement est
de 1.3~\% du spread, et pour le client, rattaché en interne à BB, industriels, États-Unis, elle prédit
337.5 points de base à dix ans (la \cref{fig:rc:the-valuation-adjustment-desk:proxy}). En appliquant
la forme de la structure par terme du chapitre~18 donne une courbe allant de 154 points de base à un an
jusqu'à 338 points de base à dix ans.
\end{example}

\begin{omfigure}
\begin{tikzpicture}
\begin{axis}[width=8.5cm,height=6cm,
  xlabel={spread observé des pairs (pb)},ylabel={spread ajusté (pb)},
  tick label style={font=\small},label style={font=\small},
  xmin=50,xmax=420,ymin=50,ymax=420,grid=major,grid style={gray!25},
  legend pos=north west,legend style={font=\footnotesize},
  table/col sep=comma]
\addplot[omExo,dashed,domain=50:420,samples=2,forget plot] {x};
\addplot[omMeth,only marks,mark=*,mark size=2pt] table[x=observed,y=fitted]{figdata/rates-credit-risk/20-the-valuation-adjustment-desk/proxy_a.csv};
\addplot[omDef,only marks,mark=square*,mark size=2pt] table[x=observed,y=fitted]{figdata/rates-credit-risk/20-the-valuation-adjustment-desk/proxy_bbb.csv};
\addplot[omThm,only marks,mark=triangle*,mark size=2.6pt] table[x=observed,y=fitted]{figdata/rates-credit-risk/20-the-valuation-adjustment-desk/proxy_bb.csv};
\legend{pairs A,pairs BBB,pairs BB}
\end{axis}
\end{tikzpicture}

\omcaption{Spreads à dix ans observés, mis en regard des spreads ajustés des pairs illustratifs de la
régression du proxy sur la notation, le secteur et la zone. Les points se situent près de la diagonale~; le
proxy du client est la prédiction du modèle pour sa case (BB, industriels, États-Unis). Données~: le
tutoriel du chapitre.}\label{fig:rc:the-valuation-adjustment-desk:proxy}
\end{omfigure}

\begin{definition}[CDS contingent]\label{def:rc:the-valuation-adjustment-desk:ccds}
Un \emph{CDS contingent}\index{CDS contingent} (contingent credit default swap) verse, lors du
défaut d'une entité de référence, la \omterm{def:rc:reduced-form-credit:pd}{perte en cas de défaut} sur la valeur alors atteinte par
un dérivé donné (ou un \omterm{def:rc:counterparty-exposure:netting}{ensemble de compensation}), et non sur un notionnel fixe~: une protection exacte pour la
CVA de ce dérivé, et tarifée au moyen de sa CVA.
\end{definition}

Une couverture par proxy laisse subsister un risque de base (le crédit du client ne suit pas celui de ses pairs),
un \omterm{def:rc:reduced-form-credit:jtd}{risque de saut au défaut} (jump-to-default)~: la protection sur indice ne verse rien quand le client fait défaut seul, et
le \omterm{def:rc:rates-risk:crossgamma}{gamma croisé} (l'exposition et le crédit évoluent ensemble). Les charges de financement et de capital n'ont
aucune couverture de marché~: le desk les gère par le collatéral, la compression, la novation et le choix des
contreparties.

\begin{dated}{2026-09}{Spreads de substitution et couvertures éligibles}\label{dat:rc:the-valuation-adjustment-desk:proxy}
Le CVA framework du Basel Committee (juillet 2020) exige que les spreads de substitution pour les
contreparties non liquides soient estimés à partir de pairs liquides par un algorithme qui discrimine au moins
selon la qualité de crédit, le secteur et la zone, autorise le rattachement à un seul nom apparenté (une municipalité
à son pays) seulement moyennant une justification prudentielle, et ne reconnaît comme couvertures de CVA que
les CDS sur nom unique, les \omterm{def:rc:the-valuation-adjustment-desk:ccds}{CDS contingents} sur nom unique et les CDS sur indice.
\end{dated}

\section{Imputation des charges, prix de cession interne et P\&L du desk}

\begin{definition}[Charge XVA]\label{def:rc:the-valuation-adjustment-desk:charge}
La \emph{charge XVA}\index{charge XVA} (\omterm{def:rc:credit-and-debit-valuation-adjustments:xva}{XVA} charge) est le montant que le \omterm{def:rc:the-valuation-adjustment-desk:desk}{desk XVA} prélève auprès du
desk d'origine lorsqu'une transaction est enregistrée~: les ajustements \omterm{def:rc:credit-and-debit-valuation-adjustments:xva}{XVA} incrémentaux de la transaction, versés d'avance ou
convertis en spread courant sur la transaction, $\text{charge} = \mathrm{XVA}/(N\cdot A)$ avec $A$ la
annuité de la transaction.
\end{definition}

\begin{example}[La cotation à vingt ans]\label{ex:rc:the-valuation-adjustment-desk:quote}
Le client reçoit le taux variable et paie le taux de pair de 4.11~\% sur 100~millions d'USD pendant
vingt ans~; sans CSA. Sur 4\,000 trajectoires, l'\omterm{def:rc:counterparty-exposure:profiles}{exposition attendue} de la banque culmine à 5.65~millions d'USD après
cinq ans (la \cref{fig:rc:the-valuation-adjustment-desk:ee}). Avec la courbe proxy, la CVA s'élève à
1\,649\,494~USD~; le financement à 80 points de base coûte 343\,872~; avec un capital imputé au \omterm{def:rc:funding-margin-and-capital-adjustments:kva}{taux de rendement minimal exigé} de 10~\%
(SA-CCR avec une pondération de risque de 100~\%, et capital CVA avec une pondération de 7~\% pour les
industriels non notés), la KVA atteint 3\,141\,826. Rapportés à l'annuité du swap de 13.67, cela fait 12.1, 2.5 et 23.0
points de base par an~: 37.6 au total.
\end{example}

\begin{omfigure}
\begin{tikzpicture}
\begin{axis}[width=11.5cm,height=5.2cm,
  xlabel={temps (années)},ylabel={EE (millions d'USD)},
  tick label style={font=\small},label style={font=\small},
  xmin=0,xmax=20,ymin=0,ymax=6.5,grid=major,grid style={gray!25},
  table/col sep=comma]
\addplot[omDef,very thick] table[x=t,y=ee]{figdata/rates-credit-risk/20-the-valuation-adjustment-desk/ee20.csv};
\end{axis}
\end{tikzpicture}

\omcaption{\omterm{def:rc:counterparty-exposure:profiles}{Exposition attendue} du swap à vingt ans, sur lequel la banque reçoit le taux fixe du
client non noté~: elle culmine après environ cinq ans et s'amortit sur les quinze années restantes.
Données~: le tutoriel du chapitre.}\label{fig:rc:the-valuation-adjustment-desk:ee}
\end{omfigure}

Le P\&L du desk est la variation des ajustements qu'il détient, nets de ses couvertures et des
charges qu'il perçoit. Les charges fixées à l'origine et les ajustements valorisés au quotidien divergent~: les spreads
évoluent, les expositions augmentent, le coût de financement propre de la banque change. Un desk qui a trop peu facturé
perd lentement~; celui qui facture trop perd les transactions.

\section{Organisation}

Le rattachement du desk détermine ce qu'il optimise. Au sein de la division marchés, il valorise et
couvre comme un desk de trading~; au sein de la trésorerie ou de la finance, il s'apparente plutôt à un service
qui fixe les prix de cession interne. Dans les deux cas, il a besoin du moteur d'exposition du chapitre~17 sur des
scénarios communs pour chaque \omterm{def:rc:counterparty-exposure:netting}{ensemble de compensation}, de données juridiques sur les accords de compensation et de collatéral, de données de crédit sur chaque
contrepartie, et d'un calculateur de capital, l'ensemble assez rapide pour répondre à un commercial en quelques secondes~:
le moteur de risque du chapitre~29.

\section{Tutoriel~: valorisation d'une nouvelle transaction}\label{tut:rc:the-valuation-adjustment-desk:tutorial}

\begin{tutorial}
\textbf{Objectif.} Calculer la CVA incrémentale et allouée dans un \omterm{def:rc:counterparty-exposure:netting}{ensemble de compensation} existant, construire un
\omterm{def:rc:the-valuation-adjustment-desk:proxy}{spread de crédit de substitution}, et coter les ajustements d'un nouveau swap de longue maturité en points de base annuels.
\textbf{État final~:} les valeurs des
\cref{ex:rc:the-valuation-adjustment-desk:incremental,ex:rc:the-valuation-adjustment-desk:euler,ex:rc:the-valuation-adjustment-desk:proxy,ex:rc:the-valuation-adjustment-desk:quote}
et les graphiques.
\begin{tutsteps}
\item \textbf{Incrémentale}~: \texttt{incremental\_table()}, autonome par rapport à l'incrémentale.
\item \textbf{Allocation}~: \texttt{allocation()}~; vérifier que les contributions se somment.
\item \textbf{Substitution}~: \texttt{proxy()} et \texttt{proxy\_curve()}.
\item \textbf{Cotation}~: \texttt{twenty\_year\_quote()}~; \texttt{fig\_rc\_xvadesk.py} produit les graphiques.
\end{tutsteps}
\textbf{Pour aller plus loin.} Rattacher le client au BBB à la place et observer le déplacement de la cotation~; ajouter une CSA
et recalculer chaque composante.
\end{tutorial}

\section{Construction~: le service de cotation}\label{bld:rc:the-valuation-adjustment-desk:xvaquote}

\begin{build}
\bfield{Rôle} Les cotations \omterm{def:rc:credit-and-debit-valuation-adjustments:xva}{XVA} de la firme~: ajustements incrémentaux d'une nouvelle transaction par rapport à
l'\omterm{def:rc:counterparty-exposure:netting}{ensemble de compensation} existant, allocations pour le reporting, spreads de substitution et charges annuelles.
\bfield{Interface} \texttt{netting\_cva(values, D, times, cpty, recovery)}~;
\texttt{incremental(adjustment, existing, new)}~; \texttt{euler\_cva}~; \texttt{proxy\_spread(peers,
target)}~; \texttt{running\_charge(upfront, notional, annuity)}. Lit \texttt{firm\_exposure},
\texttt{firm\_cva}, \texttt{firm\_xvafund}.
\bfield{Règles} Tous les ajustements portent sur les mêmes scénarios que l'\omterm{def:rc:counterparty-exposure:netting}{ensemble de compensation} existant~; les spreads de substitution
font régresser les logarithmes des spreads sur la qualité de crédit, le secteur et la région.
\bfield{Tests d'acceptation} \texttt{code/firm/xvaquote/tests/}~: les contributions d'Euler se somment et
évoluent avec la transaction~; l'identité incrémentale~; la régression du spread de substitution retrouve une
table multiplicative exactement~; la conversion de la charge annuelle.
\bfield{Pour aller plus loin} FVA et KVA incrémentales sur le portefeuille~; allocation du capital par Euler sur
la formule SA-CCR~; un spread de substitution avec des facteurs de liquidité et de séniorité~; des exécutions incrémentales
en moins d'une seconde en stockant les valeurs des trajectoires.
\end{build}

\begin{omsources}
\item Basel Committee on Banking Supervision, \emph{Targeted revisions to the credit valuation
adjustment risk framework}, juillet 2020, MAR50.
\end{omsources}

\section{Exercices}

\begin{exercise}[$\star$]\label{exo:rc:the-valuation-adjustment-desk:1}
Pourquoi la CVA incrémentale d'une transaction peut-elle être négative alors que sa CVA autonome est positive~?
\end{exercise}

\begin{exercise}[$\star$]\label{exo:rc:the-valuation-adjustment-desk:2}
Convertir une charge initiale de 1~million de dollars en points de base annuels sur un swap de 100~millions de dollars
de vingt ans avec l'annuité du chapitre.
\end{exercise}

\begin{exercise}[$\star$]\label{exo:rc:the-valuation-adjustment-desk:3}
Quelles sont les trois variables selon lesquelles un spread de substitution doit discriminer dans le cadre de la CVA
de Bâle~?
\end{exercise}

\begin{exercise}[$\star\star$]\label{exo:rc:the-valuation-adjustment-desk:4}
Montrer que les contributions d'Euler à la CVA se somment à la CVA de l'\omterm{def:rc:counterparty-exposure:netting}{ensemble de compensation}.
\end{exercise}

\begin{exercise}[$\star\star$]\label{exo:rc:the-valuation-adjustment-desk:5}
Pourquoi l'ordre d'arrivée des transactions compte-t-il pour les charges incrémentales, et pourquoi pas pour
les \omterm{def:rc:the-valuation-adjustment-desk:euler}{allocations d'Euler}~?
\end{exercise}

\begin{exercise}[$\star\star$]\label{exo:rc:the-valuation-adjustment-desk:6}
Quels risques la protection par indice laisse-t-elle ouverts lorsqu'elle couvre la CVA d'un client non noté~?
\end{exercise}

\begin{exercise}[$\star\star\star$]\label{exo:rc:the-valuation-adjustment-desk:7}
\emph{Programmation.} Calculer le \omterm{def:rc:reduced-form-credit:cs01}{CS01} de la CVA du swap de vingt ans sur la courbe de substitution et le
notionnel de protection à dix ans sur un pair qui la couvre.
\end{exercise}

\begin{exercise}[$\star\star\star$]\label{exo:rc:the-valuation-adjustment-desk:8}
\emph{Trouver la faille.} «~Notre \omterm{def:rc:the-valuation-adjustment-desk:charge}{charge XVA} couvre la perte attendue sur le client, donc le \omterm{def:rc:the-valuation-adjustment-desk:desk}{desk
XVA} devrait avoir un P\&L plat sur toute la durée de la transaction.~»
\end{exercise}

\section{Problème~: l'émetteur non noté}

\begin{problem}[{Problème du week-end --- l'ajout tout compris}]\label{pb:rc:the-valuation-adjustment-desk:1}
L'entreprise de l'\cref{ex:rc:the-valuation-adjustment-desk:quote} demande le swap de vingt ans.
Une autre banque l'a cotée 20 points de base en dessous du taux tout compris de cette banque.

\medskip\noindent\textbf{Partie I --- La cotation.}
\begin{enumerate}
\item Donner le taux au pair du swap et son annuité.
\item Donner le spread de substitution à dix ans et le classement interne qui le sous-tend.
\item Donner la CVA, la FCA et la KVA en valeur initiale et en points de base annuels.
\item Donner l'ajout tout compris.
\item Quelle composante domine, et pourquoi pour ce client~?
\end{enumerate}

\medskip\noindent\textbf{Partie II --- Couverture.}
\begin{enumerate}[resume]
\item Quelle part de l'ajout le desk peut-il couvrir sur le marché, en proportion du total~?
\item Donner le \omterm{def:rc:reduced-form-credit:cs01}{CS01} de la CVA et la couverture en protection à dix ans sur un pair.
\item Que laisse ouvert la couverture~?
\item Un \omterm{def:rc:the-valuation-adjustment-desk:ccds}{CDS contingent} serait-il disponible, et que coûterait-il~?
\item Comment le desk couvrirait-il le risque de taux de l'exposition~?
\end{enumerate}

\medskip\noindent\textbf{Partie III --- La concurrence.}
\begin{enumerate}[resume]
\item Quelles hypothèses pourraient rendre l'autre banque 20 points de base moins chère~?
\item Rattacher le client au BBB à la place~: que changerait-il~?
\item Quelles conditions la banque pourrait-elle offrir pour combler l'écart~?
\item Pourquoi l'autre banque pourrait-elle avoir tort~?
\item Renoncer est-il une perte pour la banque~?
\end{enumerate}

\medskip\noindent\textbf{Part IV --- Judgement.}
\begin{enumerate}[resume]
\item Qui doit supporter la KVA~: le desk de taux, le \omterm{def:rc:the-valuation-adjustment-desk:desk}{desk XVA} ou la banque~?
\item Comment la charge doit-elle être répartie avec le revenu du commercial~?
\item Que diriez-vous au client au sujet de la charge~?
\item Énoncer le \emph{résultat nommé}~: la surcharge tout compris en points de base courus et la part que le desk
peut couvrir.
\item En une phrase~: que vend un \omterm{def:rc:the-valuation-adjustment-desk:desk}{desk XVA} au reste de la banque~?
\end{enumerate}
\end{problem}

\section{Questions d'entretien}

\begin{interviewq}[$\star$ \iqroles{trader, banque}]\label{iq:rc:the-valuation-adjustment-desk:1}
Pourquoi les banques centralisent-elles les ajustements de valeur dans un seul desk~?
\end{interviewq}

\begin{interviewq}[$\star\star$ \iqroles{chercheur}]\label{iq:rc:the-valuation-adjustment-desk:2}
Qu'est-ce que l'\omterm{def:rc:the-valuation-adjustment-desk:euler}{allocation d'Euler}, et quand est-elle additive~?
\end{interviewq}

\begin{interviewq}[$\star\star$ \iqroles{trader}]\label{iq:rc:the-valuation-adjustment-desk:3}
Comment valoriser la CVA d'un client qui n'a aucun crédit négocié~?
\end{interviewq}

\begin{interviewq}[$\star\star$ \iqroles{développeur}]\label{iq:rc:the-valuation-adjustment-desk:4}
Un commercial a besoin d'un cours \omterm{def:rc:credit-and-debit-valuation-adjustments:xva}{XVA} incrémental en deux secondes. Comment concevriez-vous le service~?
\end{interviewq}

\begin{interviewq}[$\star\star\star$ \iqroles{trader, risque}]\label{iq:rc:the-valuation-adjustment-desk:5}
Qu'est-ce qui détermine le P\&L du \omterm{def:rc:the-valuation-adjustment-desk:desk}{desk XVA}, et que peut-il mal tourner en période de crise~?
\end{interviewq}

\begin{interviewq}[$\star\star\star$ \iqroles{banque}]\label{iq:rc:the-valuation-adjustment-desk:6}
Faut-il répercuter intégralement les \omterm{def:rc:the-valuation-adjustment-desk:charge}{charges XVA} aux clients~? Discuter la concurrence et la sélection adverse.
\end{interviewq}
